\subsection{Домашние задачи на одиннадцатую неделю.}
\begin{problem}
    Сумасшедший мирмеколог взвешивает муравьев из близлежащих муравейников.
    За летний сезон ему удается взвесить 10000 муравьев.
    За многолетний период исследований он установил, что среднеквадратическое отклонение веса муравьев за один сезон составляет 2 мг.
    В последний год исследований он обнаружил, что средний вес муравьев увеличился не менее чем на 0, 05 мг по сравнению с предыдущим годом.
    Найдите вероятность такого события.
    Не пора ли мирмекологу приобрести новые весы?
\end{problem}
\begin{solution}
    Будем считать, что выборочное среднее $X$ распределено как $N(a, \dfrac{4}{10^4})$ (согласно ЦПТ).
    Нам необходимо посчитать $\P(|X - \Ex X| > 0.05)$.
    То есть после нормировки на $\sigma$ мы получаем $\P(|Y| > \dfrac{0.05}{\sqrt {4/10^4}}) = \P(|Y| > 2.5)$, где $Y \sim N(0, 1)$.
    Тогда $\P(|Y| > 2.5) = 2\P(Y > 2.5) = 2(1 - \Phi(2.5)) \approx 0.0124$
\end{solution}
\begin{problem}
    У страховой компании $10^4$ клиентов, вероятность смерти каждого равна $6 \cdot 10^{-3}$, страховой взнос 12 у. е., выплата в случае смерти — $10^3$ у. е.
    Найти вероятность того, что доход компании превысит $4 \cdot 10^4$ у. е., а также вероятность разорения.
\end{problem}
\begin{solution}
    Клиент является бернуллиевской случайной величиной с параметром $p = 6 \cdot 10^{-3}$.
    Тогда, пусть $S = \sum\limits_{k = 1}^{10^4} Be(p)$ -- число умерших клиентов (за период наблюдения).
    Доход компании выражается как $D = 12 \cdot 10^{4} - 10^3 \cdot S$.
    Нам необходимо найти $\P(D \geq 4 \cdot 10^4)$ и $\P(D \leq 0)$, то есть $\P(S < 80)$ и $\P(S > 120)$, при $S \sim Binom(10^{4}, 6 \cdot 10^{-3})$.
    Так как $\Ex S = 60$, а $\D S = 10^{4} \cdot 6 \cdot 10^{-3} \cdot (1 - 6 \cdot 10^{-3}) = 60(1 - 6 \cdot 10^{-3})$ и по ЦПТ:
    \[
        \P\biggr(a\sqrt {\D S} + \Ex S \leq S \leq b\sqrt {\D S} + \Ex S\biggr) = \P\biggr(a \leq \dfrac{S - \Ex S}{\sqrt{\D S}} \leq b\biggr) = \Phi(b) - \Phi(a).
    \]
    Тогда в первом случае $a = -\infty, b = 2.58$ и $\P(S < 80) \approx \Phi(2.58)$, а во втором случае -- $b = +\infty, a = 7.769$, то $\P(S > 120) \approx 1 - \Phi(7.769)$.
\end{solution}
\begin{problem}
    Задачи $118$ и $121$ из канонического задачника.
\end{problem}
\begin{solution}
    В каноническом задании.
\end{solution}
\begin{problem}
    Пусть $g: \R \to \R$ -- ограниченная непрерывная функция.
    Вычислите предел $\sum\limits_{k = 0}^{+\infty} g\biggr(\dfrac{k}{n}\biggr)\dfrac{(n \lambda)^k}{k!}e^{-n\lambda}$ при $n \ra +\infty$.
\end{problem}
\begin{solution}
    Пусть $S_n \sim Poisson(n\lambda)$.
    Тогда $\Ex g(X_n/n) = \sum\limits_{k = 0}^{+\infty} g(\dfrac{k}{n})\dfrac{(n\lambda)^k}{k!}e^{-n\lambda}$ и по ЗБЧ $\dfrac{S_n}{n} \xrightarrow{\P} \lambda$, мы получаем что по лемме Портмонто $\Ex g\biggr(S_n/n\biggr) \ra \Ex g(\lambda) = g(\lambda)$.
\end{solution}
\begin{note}
    Неочевидно почему это правда и почему так можно делать, но строгое доказательство здесь проделать, кажется, достаточно нетривиально.
\end{note}
\begin{problem}
    Пусть даны $i.i.d.$ случайные величины $Y_k$, принимающие равновероятно значения $\{0, \ldots, 9\}$.
    \begin{itemize}
        \item Докажите, что $X_n = \sum\limits_{k = 1}^n \dfrac{Y_k}{10^k}$ сходятся по распределению к $U(0, 1)$.
        \item Воспользутесь тем, что $X_n$ монотонно возрастают с ростом $n$, и объясните, почему сходимость из предыдущего пункта на самом деле почти наверное.
    \end{itemize}
\end{problem}
\begin{solution} \leavevmode
    \begin{itemize}
        \item Для каждого из $Y_k$ харфункцией будет $\phi_{Y}(t) = \Ex e^{it Y} = \dfrac{1}{10}\sum\limits_{k = 0}^{9}e^{ikt} = \dfrac{1 - e^{10it}}{10(1 - e^{it})}$.
        В силу независимости величин в последовательности
        \[\phi_{X_n}(t) = \prod_{k = 1}^n \phi_Y(t/10^k) = \dfrac{1}{10^n}\prod_{k = 1}^n \dfrac{1 - e^{it \cdot 10^{1 - k}}}{1 - e^{it \cdot 10^{-k}}} = \dfrac{1 - e^{it}}{10^n(1 - e^{it \cdot 10^{-n}})} \ra \dfrac{1 - e^{it}}{-it}.\]
        Что и завершает доказательство.
        \item В силу ограниченности каждой из случайных величин и их монотонного возрастания у нас есть сходимость почти всюду.
    \end{itemize}
\end{solution}
\begin{problem}
    Пусть $X_1, \ldots$~---~$i.i.d.$ случайные величины, равномерно распределённые в $[0, 1]$.
    Вычислить среднее и дисперсию $\ln X_1$.
    Найти предел
    \[
    \lim\limits_{n \ra +\infty} \P((X_1 \cdot \ldots \cdot X_n)^{n^{-0.5}}e^{n^{0.5}} \in [a, b]).
    \]
    Применив ЦПТ к последовательности $\xi_k = \ln X_k$.
\end{problem}
\begin{solution}
   Заметим, что $-\ln X_1 \sim Exp(1)$, а следовательно $\Ex \ln X_1 = - \Ex Exp(1) = -1$ и $\D \ln X_1 = \D Exp(1) = 1$.
\end{solution}
